
\documentclass{article}
\usepackage{colortbl}
\usepackage{makecell}
\usepackage{multirow}
\usepackage{supertabular}

\begin{document}

\newcounter{utterance}

\centering \large Interaction Transcript for game `cladder', experiment `full\_v1.5\_default', episode 9908 with Qwen3.5{-}9B.
\vspace{24pt}

{ \footnotesize  \setcounter{utterance}{1}
\setlength{\tabcolsep}{0pt}
\begin{supertabular}{c@{$\;$}|p{.15\linewidth}@{}p{.15\linewidth}p{.15\linewidth}p{.15\linewidth}p{.15\linewidth}p{.15\linewidth}}
    \# & \multicolumn{2}{c}{Player} && \multicolumn{2}{c}{Game Master} \\
    \hline

    \theutterance \stepcounter{utterance}  
    & & & \multicolumn{4}{p{0.6\linewidth}}{
        \cellcolor[rgb]{0.9,0.9,0.9}{
            \makecell[{{p{\linewidth}}}]{
                \texttt{\tiny{[P1$\langle$GM]}}
                \texttt{You are an expert in causal inference. The following question is not a typical commonsense query, but rather a meticulously designed question created by a professor specializing in causal inference, intended to assess the students' mastery of the course content.} \\
\texttt{Imagine a self{-}contained, hypothetical world with only the following conditions, and without any unmentioned factors or causal relationships: Zuph has a direct effect on wibl. Wibl has a direct effect on uvzi. We know that zuph causes not wibl, and we know that wibl causes not uvzi. Would an individual is not uvzi if zuph instead of not zuph?} \\
\texttt{Start your answer with 'Yes' or 'No', followed by additional reasoning or evidence to support your explanation.} \\
            }
        }
    }
     \\ \\

    \theutterance \stepcounter{utterance}  
    & \multicolumn{4}{p{0.6\linewidth}}{
        \cellcolor[rgb]{0.561,0.737,0.561}{
            \makecell[{{p{\linewidth}}}]{
                \texttt{\tiny{[P1$\rangle$GM]}}
                \texttt{No} \\
\\ 
\texttt{**Reasoning:**} \\
\\ 
\texttt{In this hypothetical world, we must strictly adhere to the provided causal structure and the specific truth values of the relationships:} \\
\\ 
\texttt{1.  **Causal Structure**: The chain is \$Zuph \textbackslash{}rightarrow Wibl \textbackslash{}rightarrow Uvzi\$.} \\
\texttt{2.  **Specific Causal Effects**:} \\
\texttt{    *   The prompt states: "zuph causes **not** wibl." This means if \$Zuph = 1\$ (present), then \$Wibl = 0\$ (absent). Conversely, if \$Zuph = 0\$ (not present), the effect is not defined as causing "not wibl" by the specific clause, but in binary causal logic where the only alternative to "causing not X" is "causing X" (or no effect, but the prompt implies a deterministic toggle), we infer that the absence of the "not" effect implies the positive effect. Thus, if \$Zuph = 0\$, then \$Wibl = 1\$.} \\
\texttt{    *   The prompt states: "wibl causes **not** uvzi." This means if \$Wibl = 1\$, then \$Uvzi = 0\$. Conversely, if \$Wibl = 0\$, the "not" effect is absent, implying \$Uvzi = 1\$.} \\
\\ 
\texttt{3.  **Scenario Analysis**:} \\
\texttt{    *   **Case A: \$Zuph\$ is present} \\
            }
        }
    }
    & & \\ \\

    \theutterance \stepcounter{utterance}  
    & & & \multicolumn{2}{p{0.3\linewidth}}{
        \cellcolor[rgb]{0.9,0.9,0.9}{
            \makecell[{{p{\linewidth}}}]{
                \texttt{\tiny{[GM$|$GM]}}
                \texttt{No} \\
            }
        }
    }
    & & \\ \\

    \theutterance \stepcounter{utterance}  
    & & & \multicolumn{2}{p{0.3\linewidth}}{
        \cellcolor[rgb]{0.9,0.9,0.9}{
            \makecell[{{p{\linewidth}}}]{
                \texttt{\tiny{[GM$|$GM]}}
                \texttt{no} \\
            }
        }
    }
    & & \\ \\

    \theutterance \stepcounter{utterance}  
    & & & \multicolumn{2}{p{0.3\linewidth}}{
        \cellcolor[rgb]{0.9,0.9,0.9}{
            \makecell[{{p{\linewidth}}}]{
                \texttt{\tiny{[GM$|$GM]}}
                \texttt{game\_result = WIN} \\
            }
        }
    }
    & & \\ \\

\end{supertabular}
}

\end{document}
